% Placeholder template. Copy to data/cover_letter.tex and fill in your details.
%
% Header, fonts and margins deliberately match main.tex so the resume and the
% letter read as one pair of documents — keep them in step if you change either.
%
% The <<<TOKEN>>> markers are filled in by main_code/build_cover_letter.py, which
% refuses to compile if any of them survive injection. They are intentionally NOT
% LaTeX comments — a marker that leaks must be visible, not silently invisible.
\documentclass[a4paper,10pt]{article}
\usepackage[left=0.9in, right=0.9in, top=0.7in, bottom=0.7in]{geometry}
\usepackage[hidelinks]{hyperref}
\usepackage{fontspec}
\setmainfont{Times New Roman}

\pagestyle{empty}
\frenchspacing

% Block paragraphs: no indent, a clear gap between them. A letter is read as
% blocks, unlike the resume's dense sections.
\setlength{\parindent}{0pt}
\setlength{\parskip}{9pt}

\begin{document}

% Header — same identity block as the resume
\begin{center}
    {\Huge \textbf{YOUR NAME}} \\ \vspace{4pt}
    (000) 000-0000 $|$ \href{mailto:you@example.edu}{you@example.edu} $|$ \href{https://linkedin.com/in/your-handle}{linkedin.com/in/your-handle}
\end{center}

\vspace{10pt}

<<<DATE>>>

\vspace{2pt}

Hiring Team \\
<<<COMPANY>>>

\vspace{4pt}

\textbf{Re: <<<ROLE>>>}

\vspace{4pt}

Dear Hiring Team,

<<<BODY>>>

\vspace{6pt}

Sincerely, \\
Your Name

\end{document}
